% Negative fixture. Everything here is internally consistent. The suite
% asserts zero findings, which is the stronger of the two guarantees: a rule
% that fires here is a false positive, and false positives are what end this
% kind of project.
\documentclass{article}
\begin{document}

\section{Method}

Participants completed a 3-item scale adapted from prior
work~\cite{dosovitskiy2020,hu2021lora} and from~\cite{accented2019}.
We train with a learning rate of 3e-4 and batch size 128 for 100 epochs,
averaged over 5 random seeds.
% Every value agrees with the repository, and the seed is read from an
% argument, so it varies across runs. Nothing here may fire.

\section{Results}

Our approach reaches 93.8 accuracy, matching the table exactly.

\begin{table}[t]
  \caption{Accuracy and F1 on the benchmark.}
  \begin{tabular}{lcc}
    \toprule
    Method & Accuracy & F1 \\
    \midrule
    Baseline & 91.4 & 90.2 \\
    Ours     & 93.8 & 92.6 \\
    \bottomrule
  \end{tabular}
\end{table}

\begin{table}
  \caption{Dataset composition.}
  \begin{tabular}{lr}
    Split & Count \\
    Train & 800 \\
    Val   & 100 \\
    Test  & 100 \\
    Total & 1000 \\
  \end{tabular}
\end{table}

The group mean was M = 3.45 with N = 20 participants.
% Reachable: 69/20 = 3.45 exactly.

On the 3-item scale the mean was 3.47 with N = 20 participants.
% Reachable once granularity is 60 rather than 20.

The effect was reliable, $t(20) = 2.086$, $p = .05$.
The omnibus test held, $F(1, 20) = 4.35$, $p = .05$.
The association was strong, $\chi^2(1) = 6.63$, $p = .01$.

Using a one-tailed test, $t(20) = 1.725$, $p = .05$.
% Consistent one-tailed; would be inconsistent if read as two-tailed.

Across two subgroups (N = 20 and N = 25) the mean was 3.47, which cannot be
checked because the sentence names two sample sizes.

\end{document}
